% Preâmbulo comum dos diagramas (TikZ standalone). Compile com ./compilar.sh
\documentclass[tikz,border=3pt]{standalone}
\usepackage[T1]{fontenc}
\usepackage[utf8]{inputenc}
\usepackage[brazil]{babel}
\usepackage{times}
\usetikzlibrary{positioning,arrows.meta,fit,backgrounds,calc,shapes.misc}
\definecolor{azul}{HTML}{2A78D6}
\definecolor{laranja}{HTML}{EB6834}
\definecolor{verde}{HTML}{1BAF7A}
\definecolor{amarelo}{HTML}{EDA100}
\definecolor{cinza}{HTML}{52514E}
\definecolor{cinzaclaro}{HTML}{E6E5E1}
\definecolor{fundo}{HTML}{F6F5F2}
\tikzset{
  caixa/.style={draw=cinza, line width=0.5pt, rounded corners=2pt, fill=white, align=center,
                inner sep=5pt, minimum height=9mm, font=\small},
  destaque/.style={caixa, draw=azul, line width=0.8pt},
  grupo/.style={draw=cinza, dashed, rounded corners=3pt, inner sep=6pt, fill=fundo},
  rotulo/.style={font=\footnotesize\itshape, text=cinza},
  seta/.style={-{Latex[length=2mm]}, line width=0.6pt, draw=cinza},
  nota/.style={font=\scriptsize, text=cinza, align=left},
}

\begin{document}
\begin{tikzpicture}[node distance=4mm and 8mm]
  % coluna 1: o que sai de um PR real
  \node[caixa, fill=fundo, minimum width=30mm] (pr) {\textbf{PR real} do\\repositório-sujeito};
  \node[caixa, below=of pr, minimum width=30mm] (base) {Árvore no\\\emph{commit} base};
  \node[caixa, below=of base, minimum width=30mm] (enun) {Enunciado\\\scriptsize contexto, regras de negócio\\\scriptsize e contrato público};
  \node[caixa, below=of enun, minimum width=30mm, draw=laranja] (oculto) {Testes do PR\\= \textbf{oráculo oculto}};
  \draw[seta] (pr) -- (base);
  \draw[seta] (base) -- (enun);
  \draw[seta] (enun) -- (oculto);
  \begin{scope}[on background layer]
    \node[grupo, fit=(pr)(oculto), inner sep=3pt, label={[rotulo]above:uma tarefa = um PR}] (tarefa) {};
  \end{scope}
  % coluna 2: níveis
  \node[caixa, right=10mm of base, minimum width=27mm, yshift=4mm] (l0) {\texttt{none}};
  \node[caixa, below=1.5mm of l0, minimum width=27mm] (l1) {\texttt{agents}\\\scriptsize + AGENTS.md, CLAUDE.md};
  \node[caixa, below=1.5mm of l1, minimum width=27mm] (l2) {\texttt{skills}\\\scriptsize + ai/ (skills, templates)};
  \node[caixa, below=1.5mm of l2, minimum width=27mm] (l3) {\texttt{full}\\\scriptsize + opencode.json, .claude/};
  \begin{scope}[on background layer]
    \node[grupo, fit=(l0)(l3), label={[rotulo]above:4 níveis cumulativos de \emph{harness}}] (niveis) {};
  \end{scope}
  \draw[seta] (base.east) -- (niveis.west |- base.east);
  % coluna 3: agente
  \node[destaque, right=10mm of niveis, minimum width=32mm] (agente) {\textbf{Agente} (6 modelos)\\\scriptsize Claude Code ou OpenCode\\\scriptsize \emph{sandbox} Docker, só a árvore\\\scriptsize preâmbulo + enunciado\\\scriptsize até 30 min};
  \draw[seta] (niveis.east) -- (agente.west |- niveis.east);
  \draw[seta] (enun.east) -- ++(4mm,0) |- ($(niveis.south)+(0,-4mm)$) -| ($(agente.south)+(-6mm,0)$);
  % coluna 4: avaliação e saídas
  \node[caixa, right=10mm of agente, draw=laranja, minimum width=30mm] (aval) {\textbf{Avaliação}\\\scriptsize copia os testes ocultos\\\scriptsize Composer, Pint, suíte\\\scriptsize e testes ocultos isolados};
  \draw[seta] (agente) -- node[above, rotulo] {\emph{diff}} (aval);
  \node[caixa, right=7mm of aval, minimum width=28mm, yshift=6mm] (p1) {\emph{pass@1}, fração de\\testes, regressões};
  \node[caixa, below=3mm of p1, minimum width=28mm] (custo) {Custo: turnos,\\\emph{tokens}, tempo};
  \draw[seta] (aval.east) -- (p1.west |- aval.east) -- (p1.west);
  \draw[seta] (aval.east) -- (custo.west |- aval.east) -- (custo.west);
  % qualidade, sob o agente
  \node[caixa, below=14mm of agente, minimum width=60mm, draw=verde] (qual)
    {\textbf{Qualidade em 3 camadas} (sobre o \emph{diff})\\
     \scriptsize 1 Larastan, complexidade, testes escritos, Jaccard\\
     \scriptsize 2 regras C1–C10 do AGENTS.md\\
     \scriptsize 3 dois juízes LLM às cegas, rubrica de 7 dimensões};
  \draw[seta, draw=verde] (agente.south) -- node[right, rotulo, text=verde] {\emph{diff}} (qual.north);
  % oráculo oculto chega à avaliação por baixo, contornando a caixa de qualidade
  \coordinate (fundo) at ($(qual.south)+(0,-4mm)$);
  \draw[seta, draw=laranja] (oculto.east) -- ++(4mm,0) |- (fundo) -| (aval.south)
       node[pos=0.97, right, rotulo, text=laranja] {\scriptsize só na avaliação};
\end{tikzpicture}
\end{document}
